% Preface: the README's opening, its emulator, Learning Outcomes, During the Course, Getting
% Started, figures, written papers and code formatting sections; the prerequisites and the course
% plan from info/README.md; and how a lecture is arranged, from lectures/README.md, as the front
% matter of a book.

\chapter{Preface}
\markboth{Preface}{Preface}

This book is for embedded engineers who write C and have never written it for a kernel. It goes
from what a Linux system on a board is made of, through the commands that interrogate a running
one, to a device driver you write yourself: memory-mapped registers, interrupts, locking, a device
tree node, a kernel framework, and finally what all of that does to latency on a real-time kernel.

It covers Linux as embedded engineers meet it: as a kernel you configure and build, a userland you
assemble, and a driver you write for hardware nobody has written one for yet. Unlike an
introductory Linux text, it assumes you can already write C and read a datasheet, and spends its
time on the parts that are specific to kernel space:
\begin{itemize}
  \item What an embedded Linux system is made of, and the boot chain that assembles it.
  \item Open source licences, and what shipping a product built on this kernel obliges you to
    publish.
  \item The command line, taught as the set of questions you can ask a running kernel.
  \item Kernel architecture, the source tree, and Kconfig.
  \item Kernel modules: building them, loading them, and the C environment they run in.
  \item Character device drivers, and the userspace interface a driver actually presents.
  \item Kernel memory allocation, memory-mapped I/O, DMA, and managed resources.
  \item Concurrency, locking, and the contexts in which you are not allowed to sleep.
  \item Interrupts, deferred work, wait queues, and time.
  \item The driver model, the device tree, and how a driver finds its hardware.
  \item Kernel frameworks, and why a driver that is a character device is usually a driver in the
    wrong subsystem.
  \item Real-time Linux: where latency comes from, what \code{PREEMPT_RT} changes, and what makes a
    driver hostile to it.
\end{itemize}

\section*{Everything here runs in an emulator, and that is the design}
There is no board. Every chapter, every lab and every measurement runs on
\code{qemu-system-aarch64}, and the device your driver talks to is one the companion repository
ships the source of.

That is a real cost and it is worth naming before you start. You will not learn here what a
floating input feels like, what a badly seated ribbon cable does to an I2C bus, or how long a board
takes to come back after you brick its bootloader. Those are things only hardware teaches.

What it buys is everything else. The device has a register that no real device has: it records the
exact instant it raised its interrupt line, and lets your driver read that instant back. That one
register turns interrupt latency from something you infer into something you subtract, and
Chapter~\ref{c12:ch} is built on it. Beyond that, every reader gets the same machine, a wrong
\code{writel} costs you a reboot that takes two seconds rather than a trip to a bench, and a clone
of the repository plus three \code{make} targets is the whole of the setup.

\section*{What you should be able to do at the end}
Having worked through the book, you should be able to:
\begin{itemize}
  \item Name the four pieces of an embedded Linux system and say what each one does at boot.
  \item Read an open source licence well enough to say what a product built on it must publish.
  \item Configure and cross-build a kernel, and say what a given config option costs.
  \item Write, build and load an out-of-tree kernel module.
  \item Write a character device driver, and say what its \code{read()} must do that a userspace
    \code{read()} need not.
  \item Choose an allocator and a GFP flag from the context the code runs in.
  \item Map a device's registers and access them without the compiler or the CPU reordering the
    accesses out from under you.
  \item Say what address a device doing DMA actually sees, and why a buffer you may dereference is
    not automatically a buffer you may hand to hardware.
  \item Pick a lock from what it protects and from where the contention comes from, and say why
    sleeping in an interrupt handler is not a style question.
  \item Write an interrupt handler, defer the work that does not belong in it, and block a reader
    until the data exists.
  \item Write a device tree node for a device, and a driver that probes from it with no address in
    it.
  \item Say which kernel framework a device belongs to, and what registering with it gives you
    free.
  \item Explain where latency comes from on a Linux system, what \code{PREEMPT_RT} changes about
    it, and what a measurement taken under an emulator is not evidence of.
\end{itemize}

\section*{What you are assumed to know}
You are expected to arrive already comfortable with:
\begin{itemize}
  \item Programming in C at the level of writing and debugging a few hundred lines: pointers,
    structs, arrays, function pointers, and manual memory management.
  \item Bit manipulation, and reading a register map out of a datasheet.
  \item Working on a command line: files, paths, redirection, pipes, and running \code{make}.
  \item What an interrupt is, at the level of the hardware: a line, a handler, and a return.
\end{itemize}
None of that is taught from scratch. It is refreshed where a lab needs it, briefly.

What the book does teach from nothing is everything above the C: what a kernel is as a program,
what a module is as a thing you load into a running one, what a driver owes to the userspace on one
side and the hardware on the other, and why almost every rule in kernel programming turns out to be
a rule about which context the code is running in.

\textbf{The code is C, and this is not a C book.} Everything you write is ordinary C99 with a kernel
API in front of it. There are no templates, no exceptions and no standard library; that last one is
a bigger change than it sounds, and Chapter~\ref{c4:ch} spends a section on it.

\textbf{Where this book sits.} It is written to be read on its own. It assumes, and does not
derive, the process, scheduler and address-translation theory that Chapters~\ref{c3:ch}
and~\ref{c7:ch} lean on, and the memory ordering that Chapter~\ref{c6:ch} states without proving;
where a chapter needs one of them, it says so. The QAcademy course
\emph{Embedded C}\footnote{\url{https://github.com/qrtech-academy/embedded-c}} covers the C itself.

\section*{How the book is organised}
The book is typeset from a twelve-lecture course, and each chapter is one lecture. Its sections are
the lecture's appendices, in the same order and with the same subsections, so a chapter here and a
lecture directory in the repository can be read side by side. The chapters fall into four parts:

\begin{booktable}{@{}p{4.6em}p{3.6em}L@{}}
  \thead{Chapter} & \thead{Lecture} & \thead{Topic} \\
  \midrule
  \multicolumn{3}{@{}l}{\textit{Part I: The system}} \\
  1 & \file{L01} & Embedded Linux, and the licence you ship with it \\
  2 & \file{L02} & The command line, and the system underneath it \\
  3 & \file{L03} & The kernel: architecture, source tree, and configuration \\
  \addlinespace
  \multicolumn{3}{@{}l}{\textit{Part II: Writing a driver}} \\
  4 & \file{L04} & Kernel modules \\
  5 & \file{L05} & Character device drivers \\
  6 & \file{L06} & Memory, MMIO, and resource management \\
  \addlinespace
  \multicolumn{3}{@{}l}{\textit{Part III: Concurrency, interrupts and time}} \\
  7 & \file{L07} & Concurrency and locking \\
  8 & \file{L08} & Interrupts and deferred work \\
  9 & \file{L09} & Sleeping, waiting, and time \\
  \addlinespace
  \multicolumn{3}{@{}l}{\textit{Part IV: Fitting into the kernel}} \\
  10 & \file{L10} & The device model and the device tree \\
  11 & \file{L11} & Kernel frameworks \\
  12 & \file{L12} & Real-time Linux \\
\end{booktable}

Every chapter opens with what is in it, works through the material, closes with a review of what
you should now be able to explain, and ends with its exercises. The appendices hold the solutions
to the exercises that are not code, and the two written papers described below.

\textbf{Three topics are not covered here.} Yocto and Buildroot get a section in
Chapter~\ref{c1:ch} and no more, because teaching either honestly costs two chapters that would come
out of the driver material, and how to build a distribution is a different subject with different
tools. Networking drivers are absent for the same reason: the network stack is the largest
subsystem in the kernel, and a driver for it is not a first driver. CPU isolation and IRQ affinity
are absent from Chapter~\ref{c12:ch} for a different reason: the target is a two-CPU QEMU guest, and
isolating one of two CPUs from a scheduler that is itself emulated would produce numbers that say
nothing about a real machine. A claim about pinning that the target cannot demonstrate is worse
than no claim, and Chapter~\ref{c12:ch} is the chapter that argues that.

\section*{One driver, grown}
Every chapter has a lab, and from Chapter~\ref{c4:ch} onwards they are one driver, grown. You start
with a module that prints a line, and finish with a driver that probes off a device tree node,
services interrupts, exposes its device through a kernel framework, and has had its latency
measured under two different preemption models.

\textbf{The code you write is yours.} The repository ships the specification in prose, which is in
this book as the section of each chapter that specifies its lab, together with the Kbuild files,
the KUnit suites and the on-target test scripts. It never ships the driver itself.

The harness is the machine. The code you write is a kernel module, and a kernel module runs in a
kernel: there is no \code{main}, no libc, and nothing on the host that can call your
\code{read()}. So \code{make test} builds your module, drops it into an initramfs with the
lecture's \file{run.sh}, boots \code{qemu-system-aarch64}, and reads the serial console. Two things
are tested that way:
\begin{itemize}
  \item \textbf{KUnit}, for the parts that are pure logic. A ring buffer's wrap-around and a
    register field's decoding are functions with no hardware in them, and KUnit runs them inside the
    kernel with a real test framework around them. It is also what mainline uses, which makes it
    worth meeting.
  \item \textbf{The on-target script}, for everything else. Loading the module, reading \file{/dev},
    checking \code{dmesg}, counting interrupts in \file{/proc/interrupts}; all the things that are
    only true of a driver in a running kernel.
\end{itemize}
Both report through the same three markers on the console, and \code{make test} distinguishes a
test that failed from a kernel that panicked before the test began. That distinction matters more
here than it does in userspace, because in kernel space the second one is a normal Tuesday.

\section*{What the exercises look like}
Every chapter ends with eight to ten exercises, numbered within the chapter, and each labelled with
its kind:

\begin{booktable}{@{}lL@{}}
  \thead{Kind} & \thead{What it asks of you} \\
  \midrule
  \kindtag[fillaccent2]{Recall} & Say it back, in your own words, with nothing in front of you. \\
  \kindtag[fillaccent3]{Hand calculation} & Work out a number on paper. \\
  \kindtag[fillaccent4]{Design} & Choose between two mechanisms and defend the choice. \\
  \kindtag{Code} & Write something and run it. \\
  \kindtag[fillaccent]{Cross-check} & Exactly one per chapter, always last. \\
\end{booktable}

The Cross-check is the signature exercise and the reason the other four exist. You compute a number
by hand, you run your own code on the same problem, and then you reconcile the two. They will rarely
agree exactly, and reconciling them is the point: ``my two answers differ by eight percent'' is not
a failure of the exercise, it is the exercise. Each one names the artefacts to measure and the
questions to ask about the gap, so that you can tell a discrepancy you have explained from one you
have merely noticed.

Where an exercise can be checked mechanically, a \textbf{Check yourself} line says how.

\textbf{The solutions to the exercises that are not code are in Appendix~\ref{sol:ch:solutions}}:
every Recall, Hand calculation and Design exercise, and the parts of each Cross-check that are
worked on paper. \textbf{The Code exercises have no solution}, here or in the repository. Each is
checked by running it on the target, and its Check yourself line says how; an answer to read would
replace the work. Try each exercise before you turn to its solution.

\section*{The figures are drawn from code}
The nineteen figures in this book are the course's own, generated by the figure pipeline in the
repository's \file{diagrams/} directory and printed here exactly as the lectures show them. None of
them is decoration: half are block diagrams of something the prose can only describe one step at a
time, and the rest are charts of numbers the course measured on its own target. Every number on
every chart comes from one file, \file{diagrams/measured.py}, which cites the section that
publishes it, so a chart and the table beside it cannot disagree.

\section*{Two written papers, and what they are not}
Nothing in the course this book comes from is marked. Assessment is the lab after every chapter,
checked on the target by \code{make test}, and the exercises that end every chapter.

Appendices~\ref{exam:ch:about} to \ref{ansb:ch:answers} hold two four-hour papers with model
answers, and they check something else: \textbf{what you can reconstruct on paper, with nothing in
front of you.} Ten questions each, mixing prose with kernel C you either write or repair, and no
machine in the room to tell you which it is.

\textbf{They exist purely so that you can test your own knowledge after working through the book.
They gate nothing, they are not a qualification, and no part of the book requires them.}

\textbf{Take one after the last chapter.} Both papers draw on all twelve chapters, so sitting one
partway through examines material you have not read yet. Their model answers mark the code
questions against checklists rather than printing listings, so that they do not become back-door
solutions to the labs.

\section*{Setting up}
Everything builds inside a container, so the only thing you install is Docker. On Windows, install
it under WSL2 and work from the WSL side of the filesystem. Then clone the companion repository,
\repo, and from its root:

\begin{shell}
make env       # Build the container image. Once, and slow.
make kernel    # Cross-build the arm64 kernel. Once, and slower.
make qemu      # Build QEMU with the course's device in it. Once.
make rootfs    # Build the BusyBox initramfs.
make boot      # Boot the target and get a shell. Ctrl-A X quits.
\end{shell}

\noindent\code{make help} lists every target. If you would rather install the toolchain natively,
set \code{QA_NO_CONTAINER=1} and install the packages listed in \file{docker/Dockerfile}; nothing
in the course requires the container except the reproducibility of it.

Once a lab is written, \code{make build} compiles it and \code{make test} boots it and runs its
tests. Both narrow to one lecture with \code{L=}:

\begin{shell}
make build L=L06
make test L=L06
\end{shell}

\noindent Every target skips loudly rather than failing when what it needs does not exist yet, and
a skip names the target that would fix it.

\textbf{A note on the code style.} The C in the repository is formatted with \code{clang-format}
against the \file{.clang-format} shared by every QAcademy course, so the kernel code in it does
\emph{not} follow the Linux kernel's own coding style, and would be rejected upstream on that basis
alone. That is a deliberate trade for consistency with the rest of QAcademy, and
\secref{c4:app:a9} shows the same function in both styles, so that opening a file in
\file{drivers/} is not a surprise.
